\section{Closure Energy and Structural Foundations}
\label{sec:closure_energy}

\subsection{From Causal Logs to Structure}

Vol.~1 established that the causal log $C_{ij}$ serves as the fundamental
variable from which both time and space emerge:
\begin{equation}
  \frac{d\tau_i}{dt} = \rho_i \qquad \text{(time)},
\end{equation}
\begin{equation}
  d_{ij} = -\log \tilde{S}_{ij} \qquad \text{(space)},
\end{equation}
with the decomposition
\begin{equation}
  S_{ij} = \tfrac{1}{2}(C_{ij}+C_{ji}), \qquad
  A_{ij} = \tfrac{1}{2}(C_{ij}-C_{ji}).
\end{equation}
Spatial and causal structure are already encoded within the same object.
The question addressed in this volume is not how structure is generated, but:
\begin{equation}
  \boxed{\text{What configurations of this structure can persist?}}
\end{equation}

\subsection{Minimal Condition for Structural Persistence}

Consider the simplest configurations.
Two nodes define only a relation (a line); three nodes allow a loop.
A loop is the minimal condition under which a structure can reference
itself and maintain internal consistency.
\begin{equation}
  \boxed{\text{Closure first becomes possible at three nodes.}}
\end{equation}
Two-node systems cannot stabilize independently, as no enclosed relation exists.
Three-node systems introduce the first possibility of self-sustaining structure.

\subsection{Internal Representation from Logs}

Each node $i$ carries an effective internal vector $\vec{e}_i$,
defined entirely from $C_{ij}$:
\begin{equation}
  \langle \vec{e}_i, \vec{e}_k \rangle
  = \sum_{j,\ell} A_{ij} A_{k\ell} \cdot S_{ij} S_{k\ell} \cdot \tilde{S}_{j\ell}.
\end{equation}
No external embedding space is assumed.
Direction, distance, and relative orientation arise from internal
consistency conditions of the log network.
\begin{equation}
  \boxed{\text{Geometry is induced from relational structure, not presupposed.}}
\end{equation}

\subsection{Triangular Closure Measure}

For a triplet $(i,j,k)$, define the triangular closure strength:
\begin{equation}
  \phi_{ijk}
  = \exp\!\left(-\eta\,\bigl|\vec{e}_i+\vec{e}_j+\vec{e}_k\bigr|^2\right)
  \cdot \frac{1}{2}
  \left|(\vec{e}_i-\vec{e}_k)\times(\vec{e}_j-\vec{e}_k)\right|.
\end{equation}
The first factor measures closure consistency (how close the vector sum is to zero);
the second measures non-degeneracy (whether the triangle has area).
Total closure strength over a set $A$:
\begin{equation}
  \Phi(A) = \sum_{i<j<k \in A} \phi_{ijk}.
\end{equation}
For $n=2$, $\Phi = 0$; area and closure emerge first at $n=3$.

\subsection{Closure Energy Functional}

To determine which structures persist, define the closure energy:
\begin{equation}
  \label{eq:closure_energy}
  \boxed{
    E(A)
    = \alpha \left|\sum_{i\in A}\vec{e}_i\right|^2
    - \beta\,\Phi(A)
    + \gamma|A|
    + B(L(A))
  }
\end{equation}
The terms have the following roles:
\begin{center}
\begin{tabular}{ll}
\toprule
Term & Interpretation \\
\midrule
$\alpha|\sum\vec{e}_i|^2$ & penalty for incomplete closure \\
$-\beta\,\Phi(A)$ & reward for triangular closure \\
$\gamma|A|$ & structural cost \\
$B(L(A))$ & effective non-closure resistance near the complete-closure limit \\
\bottomrule
\end{tabular}
\end{center}
The final term is a phenomenological encoding of the non-closure constraint
inherited from Vol.~1.

\subsection{Particles as Stable Configurations}

Structures evolve under this energy functional.
Stable configurations correspond to local minima:
\begin{equation}
  \boxed{\text{A particle corresponds to a local minimum of } E(A).}
\end{equation}
Particles appear as persistent configurations of the underlying structure,
not as fundamental entities assumed from the outset.
\begin{equation}
  \boxed{\text{Particles are stable attractors of closure geometry.}}
\end{equation}
